\color{unimain}{\subsection{Контроль изоляции высоковольтных вводов}}
\color{black}

\begin{enumerate}

\item
Функция КИВ предназначена для диагностики состояния изоляции и защиты при возникновении повреждения высоковольтных вводов 330-750 кВ с бумажно-масляной или RIP изоляцией, устанавливаемых на автотрансформаторах и шунтирующих реакторах.

\item
Измерительные цепи защиты подключаются к потенциалметрическим выводам высоковольтных вводов с использованием согласующих трансформаторов (ТПС-0,66).

\item
Принцип действия функции КИВ основан на измерении суммы трехфазной системы токов, протекающих под воздействием рабочего напряжения через изоляцию высоковольтных вводов. При развитии повреждения изоляции высоковольтного ввода или увеличении активных утечек во время пробоя высоковольтного ввода в предварительно сбалансированной сумме токов появляется составляющая тока небаланса ($I_{\text{нб}}$) промышленной частоты, на которую реагирует функция КИВ.

\item
В нормальном режиме, даже при исправном состоянии высоковольтных вводов, сумма токов в измерительной цепи КИВ не равна нулю вследствие неравенства емкостей основной изоляции высоковольтных вводов. Для минимизации тока небаланса необходимо произвести предварительную балансировку токов фаз.

\item
Грубая балансировка токов фаз производится с помощью выбора соответствующих ответвлений первичной и вторичной обмотки согласующего трансформатора. Более точная подстройка производится путем цифрового выравнивания, используя коррекцию токов фаз по амплитуде и углу. Коррекция тока соответствующей фазы по амплитуде осуществляется путем умножения измеренного вторичного тока на соответствующий корректирующий коэффициент \textbf{«Квыр.x»}, где x – наименование фазы (A, B, C). Коррекция тока соответствующей фазы по углу осуществляется параметром \textbf{«Фкор.x»}. Расчет рабочей величины тока небаланса осуществляется после операции цифрового выравнивания путем суммирования комплексных значений токов фаз:

\begin{equation}
I_{\textsl{НБ}} = I_{\textsl{емк.А}} \cdot K_{\textsl{выр.А}} \cdot e^{\text{jФкор.A}} + I_{\textsl{емк.B}} \cdot K_{\textsl{выр.B}} \cdot e^{\text{jФкор.B}} + I_{\textsl{емк.C}} \cdot K_{\textsl{выр.C}} \cdot e^{\text{jФкор.C}}
\end{equation}

\begin{eqexpl}[10mm]
\eqexplSetItemAlign{r}  % выравнивание обозначений по левому краю(l)/правому краю(r)/центру(c)
\item{$I_{\textsl{емк.x}}$} комплексное значение вторичного емкостного тока утечки в соответствующей фазе;
\item{$K_{\textsl{выр.x}}$} коэффициент амплитудного выравнивания (индивидуально для каждой фазы);
\item{Ф$_{\textsl{кор.x}}$} значение коррекции по углу в градусах (индивидуально для каждой фазы).
\end{eqexpl}

\item
Функция КИВ содержит сигнальный, отключающий и блокирующий органы, а также избиратели поврежденной фазы.

\item
Сигнальный орган указывает на начавшееся прогрессирующее повреждение изоляции высоковольтного ввода. Значение срабатывания сигнального органа задается параметром \textbf{«Iсигн»}. Сигнальный орган с задержкой \textbf{«Тсигн»} действует на предупредительную сигнализацию.

\item
При срабатывании сигнального органа осуществляется индикация повреждённой фазы высоковольтного ввода. Принцип работы избирателя повреждённой фазы основан на сравнении модуля емкостного тока каждого из вводов с модулем геометрической суммы емкостных токов вводов двух других фаз:

\begin{equation}
\lvert I_{\textsl{емк.A}}' \rvert \geq \lvert I_{\textsl{емк.B}}' + I_{\textsl{емк.C}}'\rvert, \lvert I_{\textsl{емк.B}}' \rvert \geq \lvert I_{\textsl{емк.A}}' + I_{\textsl{емк.C}}'\rvert, \lvert I_{\textsl{емк.C}}' \rvert \geq \lvert I_{\textsl{емк.A}}' + I_{\textsl{емк.B}}'\rvert
\end{equation}

\begin{eqexpl}[10mm]
\eqexplSetItemAlign{r}  % выравнивание обозначений по левому краю(l)/правому краю(r)/центру(c)
\item{$I_{\textsl{емк.x}}'$} комплексное значение вторичного емкостного тока утечки в соответствующей фазе после операции цифрового выравнивания.
\end{eqexpl}

\item
Отключающий орган действует при дальнейшем развитии повреждения высоковольтных вводов. Значение срабатывания отключающего органа задается параметром \textbf{«Iоткл»}. Отключающий орган с задержкой \textbf{«Тоткл»} действует на отключение защищаемого оборудования.

\item
Для исключения ложной работы сигнального и отключающего органа при внешних КЗ на землю в сети ВН и при неполнофазных режимах сети контролируется значение напряжения 3U0. При превышении уровнем напряжения 3U0 значения \textbf{«3U0макс»} производится загрубление сигнального и отключающего органов. Также загрубление может быть произведено с использованием внешнего дискретного сигнала. В режиме загрубления пороги срабатывания сигнального и отключающего органов определяются параметрами \textbf{«Iсигн.груб»} и \textbf{«Iоткл.груб»} соответственно, а отключение защищаемого оборудования производится с дополнительной задержкой \textbf{«Тдоп.груб»}.

\item
Блокирующий орган предназначен для исключения ложной работы отключающего органа при обрывах или нарушении изоляции кабеля между измерительными выводами и согласующим трансформатором. Порог срабатывания отключающего органа определяется параметром \textbf{«Iблк»}. При срабатывании блокирующего органа с задержкой по времени \textbf{«Тсигн.неисп»} осуществляется индикация неисправной фазы высоковольтного ввода. Выбор неисправной фазы осуществляется по снижению в ней уровня тока ниже значения \textbf{«Iнеисп»}.

\item
Пороги срабатывания задаются относительно базисного тока, значение которого определяется по формуле:

\begin{equation}
I_{\textsl{баз}} = \frac{I_{\textsl{ном.A}} \cdot I_{\textsl{ТПС.A}} + I_{\textsl{ном.B}} \cdot I_{\textsl{ТПС.B}} + I_{\textsl{ном.C}} \cdot I_{\textsl{ТПС.C}}}{3}
\end{equation}

\begin{eqexpl}[10mm]
\eqexplSetItemAlign{r}  % выравнивание обозначений по левому краю(l)/правому краю(r)/центру(c)
\item{$K_{\textsl{ТПС}}$} коэффициент трансформации согласующего трансформатора соответствующей фазы;
\item{$I_{\textsl{ном}}$} номинальный первичный емкостной ток утечки в соответствующей фазе, А.
\end{eqexpl}

\begin{equation}
I_{\textsl{ном}} = \frac{314 \cdot U_{\textsl{ном}} \cdot C_{\textsl{ном}}}{\sqrt{3}}
\end{equation}

\begin{eqexpl}[10mm]
\eqexplSetItemAlign{r}  % выравнивание обозначений по левому краю(l)/правому краю(r)/центру(c)
\item{$U_{\textsl{ном}}$} номинальное линейное напряжение, В;
\item{$C_{\textsl{ном}}$} паспортное значение емкости основной изоляции ввода в соответствующей фазе, Ф.
\end{eqexpl}

\item 
Уставки КИВ приведены в таблице \ref{app_set:kiv}. Алгоритм работы КИВ приведен на рисунке \ref{pict_kiv:KIV}.

\medskip
\begin{figure}[H]
\centering
\includegraphics[width=1\textwidth,height=1\textheight,keepaspectratio]{\fbpath/11.01 ШР/161.1131 КИВ/img/KIV.pdf}
\captionof{figure}{Функциональная схема логики КИВ}\label{pict_kiv:KIV}
\end{figure}


\end{enumerate}